\chapter*{}
\addcontentsline{toc}{chapter}{Introduction Générale}
\markboth{Introduction Générale}{Introduction Générale}

\begin{center}
{\fontsize{20}{24}\selectfont\bfseries\color{black} Introduction Générale}
\end{center}

\vspace{12pt}

Dans un environnement économique marqué par une concurrence accrue et des mutations 
technologiques profondes, la maîtrise des opérations commerciales s'impose comme un levier 
stratégique incontournable. Les dépôts de vente de boissons en sont une illustration concrète : 
assurer un suivi rigoureux de l'activité commerciale et piloter efficacement les performances 
au quotidien nécessite des outils pensés pour leur réalité métier. Or, à mesure que les volumes 
de données s'accumulent, savoir les exploiter à des fins décisionnelles cesse d'être un luxe 
pour devenir un avantage concurrentiel déterminant, car la performance ne repose pas uniquement 
sur la qualité des produits, mais sur la faculté de l'entreprise à transformer ses données en 
décisions éclairées.

Les approches traditionnelles de gestion commerciale, souvent fondées sur des outils numériques 
non intégrés, des tableaux de bord génériques et des fonctionnalités inadaptées aux besoins 
spécifiques de chaque activité, génèrent des inefficacités considérables. Les responsables 
peinent à obtenir une vision globale et en temps réel de leur activité commerciale. Les 
encaissements sont difficiles à suivre, les tendances de vente restent inexploitées, et 
l'absence de mécanismes de prévision contraint les décideurs à s'appuyer sur des estimations 
empiriques plutôt que sur des données structurées. Ces lacunes freinent la réactivité 
opérationnelle et limitent considérablement la qualité du pilotage commercial.

Pour répondre à ces défis, notre projet vise à concevoir et développer un système complet de 
gestion et de suivi des ventes, réalisé au sein de Teknovat Solutions, une start-up tunisienne 
spécialisée dans le développement de solutions logicielles sur mesure, pour le compte d'un 
dépôt de boissons (eaux, jus, boissons gazeuses) dont la clientèle est principalement composée 
de cafés et de restaurants. La solution proposée intègre sept modules fonctionnels couvrant la 
gestion des agents, des produits, des clients, des ventes, des factures, des règlements et des 
notifications, dont l'envoi de factures par WhatsApp. Elle est enrichie par des tableaux de 
bord interactifs dédiés respectivement au chiffre d'affaires, à la facturation et aux règlements, 
avec des indicateurs clés de performance calculés automatiquement, ainsi qu'un module de 
prévision des ventes basé sur la régression linéaire afin d'anticiper les évolutions futures 
de l'activité.

Le développement de notre système suit le cadre du Processus Unifié (UP), une méthodologie 
structurée et itérative garantissant une progression rigoureuse depuis la collecte des besoins 
jusqu'à la mise en œuvre finale. Cette approche s'appuie sur le langage de modélisation UML 
pour la conception du système, combinée à la méthode GIMSI pour la construction de tableaux 
de bord décisionnels alignés sur les objectifs stratégiques de l'entreprise.

Ce rapport de projet de fin d'études est organisé en quatre chapitres :

\begin{itemize}
    \item \textbf{Chapitre 1 --- Contexte général :} présente l'organisme d'accueil 
    Teknovat Solutions, analyse le système existant et ses limites, et expose les objectifs 
    du projet, la solution proposée ainsi que la méthodologie adoptée.

    \item \textbf{Chapitre 2 --- Expression des besoins :} identifie les acteurs du système 
    et leurs interactions, détaille les besoins fonctionnels, non fonctionnels et décisionnels, 
    et présente la spécification complète des cas d'utilisation.

    \item \textbf{Chapitre 3 --- Conception :} décrit l'architecture logique du système, la 
    conception détaillée de chaque cas d'utilisation, le schéma de base de données, ainsi que 
    la conception des tableaux de bord et du module de prévision.

    \item \textbf{Chapitre 4 --- Implémentation :} présente l'environnement logiciel, les 
    technologies utilisées et les interfaces finales du système, illustrant la concrétisation 
    de la solution développée.
\end{itemize}